\label{sec:related}

The literature most relevant to this manuscript falls into seven closely connected
groups: high-level finite-element automation, PDE-level adjoint optimization,
JAX-native differentiable FEM, JAX-to-FEM/PETSc bridge mechanisms,
Mohr--Coulomb slope-stability source-lineage context, topology optimization
workflows, and sparse nonlinear solver infrastructure. The purpose of this
section is not to assign a single ranking to those groups, but to state precisely
which part of each literature thread the present paper uses as context.

\paragraph{High-level finite-element automation.}
The original \fenics{} project established the high-level finite-element
programming model in which weak forms are expressed symbolically and compiled
into efficient kernels \citep{logg2012fenicsbook}. The more recent
\code{DOLFINx} preprint keeps that high-level mathematical abstraction while
moving to a data-oriented design with explicit support for modern parallelism,
custom kernels, and direct access to lower-level data structures
\citep{baratta2025dolfinx}. This literature defines the reference standard for
what an expressive open-source FEM environment should provide.

\paragraph{Adjoint-based PDE optimization.}
\code{dolfin-adjoint} showed that adjoints of high-level transient finite-element
programs can be derived automatically from the symbolic problem description
\citep{farrell2013dolfinadjoint}. The later \code{pyadjoint} software paper
documents the maintained FEniCS/Firedrake adjoint framework used to automate
such workflows in practice \citep{mitusch2019pyadjoint}. The present
\code{cashocs} software paper extends that PDE-constrained-optimization thread
to shape optimization, optimal control, topology optimization via level sets,
and MPI-enabled workflows \citep{blauth2023cashocsv2}. These works are the
right comparison point for adjoint-centered PDE-level automation and
optimization workflows, but they do not have the same focus as the present
paper on explicit second-order routes inside distributed sparse nonlinear solve
stacks.

\paragraph{JAX-native differentiable FEM.}
The broad AD background is surveyed by \citet{baydin2018autodiff}, while JAX
provides the program-transformation environment that makes higher-order
derivatives, vectorization, and just-in-time compilation readily usable in
scientific computing \citep{jax2018}. Within finite elements, JAX-FEM shows that
inverse design, nonlinear constitutive modeling, and GPU-oriented differentiable
workflows can be implemented inside a JAX-native 3D FE solver
\citep{xue2023jaxfem}. Xue's later finite-element differentiable-physics paper
is a direct second-order comparator because it derives implicit Hessian-vector
products and evaluates Newton-CG-style use of exact Hessian information on
finite-element inverse-problem benchmarks \citep{xue2026implicit}. AutoPDEx
broadens that JAX-native PDE picture with nonlinear minimizers, implicit
differentiation, and optional PETSc integration \citep{bode2025autopdex}.
JAX-CPFEM is especially relevant to the mechanics side of the present paper
because it couples differentiable FEM, nonlinear crystal plasticity, and GPU
execution in a recent open-access workflow \citep{hu2025jaxcpfem}. Together,
these papers define the closest modern literature context for the repository's
JAX-based local differentiation layer and second-order derivative comparisons.

\paragraph{Bridges between JAX-style AD and established FEM stacks.}
The Firedrake--JAX bridge of \citet{yashchuk2023bringing} is important because
it bypasses differentiation through low-level solver iterations by using
tangent-linear and adjoint equations derived from the high-level PDE
representation. The recent \code{FEniCSx} external-operator paper by
\citet{latyshev2025externaloperators} is even closer to the constitutive story
of this manuscript: it uses algorithmic automatic differentiation for general
constitutive models inside \code{FEniCSx}, including a Mohr--Coulomb example.
Those papers show that constitutive- or solve-level differentiation can be
combined with established FEM environments without differentiating every low-level
algebraic step. The present paper extends that bridge-style differentiation
story into a PETSc-centered sparse assembly and nonlinear-solve setting with
explicit comparison between element AD, constitutive AD, and colored SFD. The
JetSCI preprint is the closest recent hybrid in architecture: it uses JAX for
local finite-element discretization and PETSc for robust distributed sparse
solves in heterogeneous micromechanics \citep{cattaneo2026jetsci}. We therefore
treat it as evidence from a submitted preprint that the JAX-local/PETSc-global
split is an active design point, while keeping the present paper's claim
narrower: a maintained nonlinear energy-minimization suite with derivative-route
comparisons across the benchmark families below.

\paragraph{Mohr--Coulomb slope stability and source-lineage context.}
The plasticity benchmarks use Mohr--Coulomb-type slope-stability problems as
reviewer-visible stress tests, so their citations have to distinguish continuum
plasticity, strength-reduction practice, and repository-specific surrogates.
Davis' plasticity treatment and the finite-element strength-reduction studies of
\citet{tschuchnigg2015nonassociated} provide the Davis-style reduction context
used in the paper. For constitutive integration, the incremental reference frame
is return mapping with consistent tangents \citep{simo1985consistent,simo1998compinel},
with Mohr--Coulomb-specific nonsmooth return operators treated by
\citet{sysala2017returnmapping}. The recent Sysala slope-stability line also
documents variational and optimization viewpoints for shear strength reduction
\citep{sysala2021optimization,sysala2025convexoptimization} and a 3D
continuation/iterative-solver source family \citep{sysala2025advancedcontinuation}.
The present manuscript uses that literature to position the benchmark family,
while the numerical claims remain limited to the implemented surrogate and the
repository artifacts reported below.

\paragraph{Topology optimization workflows.}
The classical SIMP and educational topology-optimization references remain
important because they supply the baseline compliance, filtering, and
continuation ideas used throughout the field
\citep{sigmund2001topology,bendsoe2003topology,ferrari2020top99,bourdin2001filters}.
For the present paper, however, the more direct modern comparator is
\code{FEniTop}, which documents a simple 2D/3D \code{FEniCSx} topology-optimization
implementation with parallel support \citep{jia2024fenitop}. That paper is
relevant not because it duplicates our reduced design objective, which it does
not, but because it provides a recent source-backed baseline for how a compact
parallel topology-optimization code can be organized in the modern FEniCSx
ecosystem.

\paragraph{Sparse derivative recovery and scalable nonlinear infrastructure.}
When exact second-order information is too expensive, sparse derivative recovery
via graph coloring remains the classical alternative. Curtis, Powell, and Reid
gave the foundational sparse-Jacobian viewpoint \citep{curtis1974sparsejacobian},
while Coleman and Mor{\'e} extended the approach to sparse Jacobian and Hessian
estimation with graph-coloring structure
\citep{coleman1983sparsejacobian,coleman1984sparsehessian}. PETSc supplies the
distributed vectors, matrices, Krylov methods, nonlinear solvers, and multigrid
infrastructure needed to turn those derivative objects into scalable workflows
\citep{balay1997petsc,petsc2024web}. In this paper, PETSc is not treated as a
mere back-end library: ownership layout, Krylov policy, nonlinear globalization,
and preconditioner design are all part of the scientific comparison.

Against that background, the niche of \repo{} can be stated precisely. It sits
between high-level differentiable modeling and scalable sparse nonlinear solver
engineering. The contribution is not that it replaces the ecosystems above, but
that it documents one specific point in the contemporary design space:
maintained distributed nonlinear FEM workflows that combine JAX-based local
differentiation, multiple second-order routes, and PETSc-based sparse execution
on benchmark families ranging from scalar nonlinear diffusion to topology
optimization and 3D Mohr--Coulomb-type mechanics.
